% !TeX root = ../../main.tex

\section{結論}\label{sec:結論}

綜合教學現場調查與文獻回顧，我們發現：

\subsection{(一) AI帶來哪些教育現場的效益？}\label{sec:結論一}

\begin{itemize}
    \item \textbf{學生端}：提供更具個別化的學習環境，如個別化回饋、解決能力差異的配對問題、個別化的指導等。正因為學生能夠對AI提出任何提問，而AI也能夠給出回覆，不僅能夠對學生給予個別化的回饋，讓學生能夠在任何時間點提出任何問題，且重複提問直到真正理解為止，也能夠讓AI換一種方式來說明，以配合學生的理解程度差異，這樣的學習支持有效的改善了過去教育科技的侷限，過去的教育科技多侷限於某個學科或單元，且使用方式較為侷限，學生可操作的空間較小，而生成式AI相較於其他教育科技具有更大的彈性，讓學生可以根據自己所需來使用AI工具，為個別化學習提供更大的可能性。

    \item \textbf{教師端}：AI幫助教師在備課或作業批改上有更多協助，在教學現場端有許多重複性的工作，如作業批改、成績計算等，若AI能夠幫助教師完成這些部分，教師則能夠更專注於教學設計或提供學生更個別化的學習建議，以減輕教師的負擔。
\end{itemize}

此外，相關研究指出AI能夠帶來提升學習成效、降低焦慮、幫助概念理解、降低認知負荷、提升認知參與度、提升自我效能等幫助，顯示了其在各個方面為學習帶來的助益。

\subsection{(二) AI為教育帶來哪些擔憂？}\label{sec:結論二}

從現場調查與研究回顧可見，許多教師擔憂學生使用AI產生的依賴問題，即便學生透過AI可以完成更多的學習任務，並減輕學習壓力，然而許多研究指出，在AI融入教育的情況下，更應注重學生\textbf{批判性思考能力}，因學生要能夠判斷AI提供的內容是否正確，就必須具備批判AI提供的內容的能力，而非全然相信，尤其針對此種非針對教育而設計的大型語言模型，產出的內容往往可能有誤，若學生不具備批判性思考的能力，則最終可能學習到錯誤的知識。此外，在使用方法上，若學生僅向AI索取答案，跳過了思考的部分，則會讓過去的評量標準失去其功能性，即便學生繳交出正確答案，仍然無法確保其真正理解，因此在AI融入教學的設計中，其評量方式需要做出改變。

\subsection{(三) AI是否影響教育的目標？}\label{sec:結論三}

十二年國教課程綱要總綱中指出：

\begin{quote}
「核心素養」是指一個人為適應現在生活及面對未來挑戰，所應具備的知識、能力與態度。

「國家課程規範歷經數次中小學課程標準修訂，務求課程修訂能與時俱進。以培養健全國民為宗旨，為我國人才培育奠定良好基礎。」
\end{quote}

課綱的修訂會隨著時代變遷而有所修訂，目前AI已經取代部分工作時，顯示了AI已開始對社會造成改變，人們已經不可避免的會接觸並使用AI，因此教育部在未來的課綱修訂中，應考量如何使學生具備使用AI的能力，並且同時兼具在「自主行動」、「溝通互動」、「社會參與」等面向的學習目標，AI融入皆可能會對這些面向產生衝擊，例如：在「自主行動」面向中

\begin{quote}
「應培養學生具備問題理解、思辨分析、推理批判的系統思考與後設思考素養，並能行動與反思，有效處理及解決生活、生命問題。」
\end{quote}

若使用AI時直接獲取答案，則無法培養問題理解、思辨分析、推理批判等能力，因此未來在AI融入教學上應謹慎評估學生是否具備這些能力，且應培養學生對AI的了解，以及使用AI時具備批判性思考的能力。

\subsection{(四) AI會做這些事情，我們還要學習嗎？}\label{sec:結論四}

\begin{quote}
「教育是為適應現在生活及面對未來挑戰，成為具有社會適應力與應變力的終身學習者。」
\end{quote}

AI幫助我們的生活更便利，讓我們可以不用付出相等的努力就完成過去無法完成的事情，然而這不代表著我們就能夠因此不去學習，而是應該從新思考在AI時代下，我們應該學習的事情是什麼。隨著AI快速發展與變遷，有必要不斷地檢核教育的目標與方法，以適應AI以對社會帶來的改變。人的時間及力氣是有限的，應該進一步思考在未來的課綱中那些目標應該調整或加入，高階能力的培養起始於基礎能力，因此不應因為AI能夠取代這些基礎能力，就認為培養這些基礎能力並不重要，可以重新調整這些課程設計上的比重，以培養學生在各項指標上的能力。

\section{未來建議}\label{sec:未來建議}

\subsection{(一) 課綱的調整}\label{sec:未來建議一}

面對AI時代，十二年國教課綱應重新思考並納入AI作為一門學科或是AI融入教學的設計，並將其納入電腦課程或是強化其他課程在批判性思考的培養，以健全學生在面對AI變遷的影響，並考量到AI普及化下教育目標應如何調整，確保學生有足夠的能力來應對社會發展。

\subsection{(二) AI融入教學設計}\label{sec:未來建議二}

在AI融入教案的設計上，應重新去檢視AI如何融入教學，以及過去評量方法是否還能夠有效評估學生的能力，在教師端也應特別去設計能夠檢核學生是否真正學會的環節，可透過讓學生紀錄AI使用的方法、批判性思考環節等，避免學生過度依賴AI的問題。

\subsection{(三) 國家的AI使用指引}\label{sec:未來建議三}

面對生成式 AI 對教育體系的衝擊，臺灣亟需整合國際政策經驗，制定具體且具可操作性的國家級使用指引。本報告結合跨國政策比較與本土調查，提出以下五大核心政策建議：

第一，\textbf{建立分齡分層的 AI 素養教育體系}。我國應以教育部現行的《我和 AI 一起學》與《駕馭 AI，洞察未來》手冊為基礎\cite{govtw2024aiguide}，結合聯合國教科文組織（UNESCO）的「理解、應用、創造」三階段漸進模式\cite{unesco2024competency}，建構系統化素養指標。國小階段著重於辨識 AI 幻覺與隱私保護；國高中階段強調批判性查核與共學；高等教育則專注於高階提示詞工程與倫理責任。

第二，\textbf{重構評量機制，轉向過程驗證}。為避免如美國普渡大學因 AI 偵測軟體誤判而引發的校園信任危機\cite{arun2026ai}，評量機制必須從「結果導向」轉為「過程導向」。各級學校應減少回家作業的分數比重，並效法全國中小學科展的口頭答辯機制\cite{nehs2025science}，透過即席問答與實作演示，確認學生真正內化知識，而非僅是複製 AI 的產出。

第三，\textbf{落實強制揭露與使用者追責制度}。在學術與競賽中，應確立「創意歸屬人類，勞務外包 AI」的原則。參考新加坡國立大學的當責要求\cite{nus2024integrity}與旺宏科學獎的作法\cite{mxic2025competition}，強制要求學生在提交作業時，必須完整保留並揭露與 AI 協作的提示詞軌跡（Prompt Logs），且使用者須對所有產出內容負起最終責任。

第四，\textbf{推動教師賦能與系統性培訓}。借鑑英國教育部的雙軌策略\cite{dfe2023genai}與新加坡高達百分之七十五的教師 AI 使用率\cite{moe2025singapore}，國家應提供系統性的教師培訓，協助教師使用 AI 工具減輕行政與備課負擔，並提升教師設計 AI 融入課程與評估學生 AI 使用軌跡的專業能力。

第五，\textbf{強化風險評估與資料隱私保護}。在推動 AI 融入教學的同時，必須建立嚴格的安全防護網。參考英國教育部對網路監控與提示詞可見性的要求\cite{dfe2023genai}，以及日本文部科學省對學生思考力低下的警示\cite{mext2024genai}，學校在導入 AI 工具前應進行全面的風險評估，確保學生隱私安全，並引導學生維持主體思考，避免盲目接受 AI 答案。

\section{組員分工}\label{sec:組員分工}

\begin{longtable}{ll}
\toprule
成員 & 分工 \\
\midrule
陳慶霖 & 摘要、問卷設計與發放、文獻回顧與分析、結論、未來建議 \\
鍾詠傑 & 議題簡介、現況分析、未來建議 \\
林吟貞 & 學生態度分析 \\
林駿甫 & 教師觀點分析 \\
\bottomrule
\end{longtable}

\addcontentsline{toc}{section}{參考文獻}
\printbibliography[segment=\the\value{refsegment},title=參考文獻,heading=subbibliography]
